\subsection{模的森田对应}

在本小节中，字母 A 与 B 表示森田等价的 k-代数，P 表示一个可逆 \((A,B)_{k}\)-双模。取定一个与 P 互逆的 \((B,A)_{k}\)-双模 Q 以及同构

\[
\lambda : \mathrm{P} \otimes_ {\mathrm{B}} \mathrm{Q} \rightarrow {} _ {s} \mathrm{A} _ {d} \quad \text { and } \quad \theta : \mathrm{Q} \otimes_ {\mathrm{A}} \mathrm{P} \rightarrow {} _ {s} \mathrm{B} _ {d}.
\]

对每个左 B-模 V，用 \(\theta_{V}\) 表示 B-模同态 \(\theta\otimes1_{V}:Q\otimes_{A}P\otimes_{B}V\to V\)；由于 \(\theta\) 是同构，它是同构。同样，对每个左 A-模 M，用 \(\lambda_{M}\) 表示 A-模同态 \(\lambda\otimes1_{M}:P\otimes_{B}Q\otimes_{A}M\to M\)；由于 \(\lambda\) 是同构，它是同构。

定理 2（森田）. —— a) 设 V 与 W 是左 B-模。映射 \(g \mapsto 1_{\mathrm{P}} \otimes g\) 是从 \(\operatorname{Hom}_{\mathrm{B}}(\mathrm{V}, \mathrm{W})\) 到 \(\operatorname{Hom}_{\mathrm{A}}(\mathrm{P} \otimes_{\mathrm{B}} \mathrm{V}, \mathrm{P} \otimes_{\mathrm{B}} \mathrm{W})\) 的双射。逆双射把 \(\operatorname{Hom}_{\mathrm{A}}(\mathrm{P} \otimes_{\mathrm{B}} \mathrm{V}, \mathrm{P} \otimes_{\mathrm{B}} \mathrm{W})\) 的元素 h 映为元素 \(\theta_{\mathrm{W}} \circ (1_{\mathrm{Q}} \otimes h) \circ \theta_{\mathrm{V}}^{-1}\)，它属于 \(\operatorname{Hom}_{\mathrm{B}}(\mathrm{V}, \mathrm{W})\)。

\begin{enumerate}
\def\labelenumi{\alph{enumi})}
\setcounter{enumi}{1}
\tightlist
\item
  每个左 A-模 M 都同构于形如 \(\mathrm{P} \otimes_{\mathrm{B}} \mathrm{V}\) 的模，其中 V 是左 B-模。
\end{enumerate}

设 V 与 W 是左 B-模。由 VIII, 第 101 页的引理 1，映射 \(\varphi : g \mapsto 1_{\mathrm{P}} \otimes g\)（从 \(\mathrm{Hom}_{\mathrm{B}}(\mathrm{V}, \mathrm{W})\) 到 \(\mathrm{Hom}_{\mathrm{A}}(\mathrm{P} \otimes_{\mathrm{B}} \mathrm{V}, \mathrm{P} \otimes_{\mathrm{B}} \mathrm{W})\)）是单射。交换 P 与 Q（以及 A 与 B）的角色，可见映射 \(\psi : h \mapsto 1_{\mathrm{Q}} \otimes h\)（从 \(\mathrm{Hom}_{\mathrm{A}}(\mathrm{P} \otimes_{\mathrm{B}} \mathrm{V}, \mathrm{P} \otimes_{\mathrm{B}} \mathrm{W})\) 到 \(\mathrm{Hom}_{\mathrm{A}}(\mathrm{Q} \otimes_{\mathrm{A}} \mathrm{P} \otimes_{\mathrm{B}} \mathrm{V}, \mathrm{Q} \otimes_{\mathrm{A}} \mathrm{P} \otimes_{\mathrm{B}} \mathrm{W})\)）也是单射。而复合 \(\psi \circ \varphi\) 是映射 \(g \mapsto \theta_{\mathrm{W}}^{-1} \circ g \circ \theta_{\mathrm{V}}\)。它是双射，因而 \(\psi\) 也是。于是 \(\varphi\) 是双射，其逆是映射 \(h \mapsto \theta_{\mathrm{W}} \circ (1_{\mathrm{Q}} \otimes h) \circ \theta_{\mathrm{V}}^{-1}\)。

断言 b) 由 \(\lambda_{\mathrm{M}}\) 是从 \(\mathrm{P} \otimes_{\mathrm{B}} \otimes_{\mathrm{A}} \otimes_{\mathrm{M}}\) 到 \(\mathrm{M}\) 的同构这一事实得出。

设 V 是一个左 B-模，W 是 V 的一个子模。由于 B-模 P 是投射的（VIII, 第 101 页，定理 1），从 \(P \otimes_{B} W\) 到 \(P \otimes_{B} V\) 的典范映射是单射。我们借此映射把 \(P \otimes_{B} W\) 与它在 \(P \otimes_{B} V\) 中的像等同起来。当 P 与 B 分别换成 Q 与 A 时，我们采用类似的约定。

命题 5. —— 设 V 是一个左 B-模。映射 \(W \mapsto P \otimes_{B} W\) 是从 V 的 B-子模之集（按包含关系排序）到 \(P \otimes_{B} V\) 的 A-子模之集（按包含关系排序）的同构。逆同构把 \(P \otimes_{B} V\) 的 A-子模 N 映为 \(\theta_{V}\) 作用下 B-子模 \(Q \otimes_{A} N\)（\(Q \otimes_{A} P \otimes_{B} V\) 的子模）的像。

用 \(D_{\mathrm{B}}(\mathrm{V})\) 表示 V 的按包含关系排序的 B-子模之集，并同样地定义集合 \(D_{\mathrm{A}}(\mathrm{P}\otimes_{\mathrm{B}}\mathrm{V})\) 与 \(D_{\mathrm{B}}(\mathrm{Q}\otimes_{\mathrm{A}}\mathrm{P}\otimes_{\mathrm{B}}\mathrm{V})\)。设 \(\varphi:D_{\mathrm{B}}(\mathrm{V})\to D_{\mathrm{A}}(\mathrm{P}\otimes_{\mathrm{B}}\mathrm{V})\) 是映射 \(W\mapsto P\otimes_{B}W\)，\(\psi\) 是从 \(D_{\mathrm{A}}(\mathrm{P}\otimes_{\mathrm{B}}\mathrm{V})\) 到 \(D_{\mathrm{B}}(\mathrm{Q}\otimes_{\mathrm{A}}\mathrm{P}\otimes_{\mathrm{B}}\mathrm{V})\) 的映射，由 \(N\mapsto Q\otimes_{A}N\) 给出。它们都是递增映射，且复合 \(\psi\circ\varphi\) 是映射 \(W\mapsto\theta_{\mathrm{V}}^{-1}(\mathrm{W})\)，后者是双射。于是 \(\varphi\) 是单射，而 \(\psi\) 是满射。把 B 换成 A、V 换成 \(P\otimes_{B}V\)，即见 \(\psi\) 也是单射。从而 \(\varphi\) 与 \(\psi\) 都是双射，且 \(\varphi\) 的逆映射正是命题中所描述的那个。

例. —— 1) 把 VIII, 第 104 页的命题 5 应用于特殊情形 \(\mathrm{V} = \mathrm{B}_s\)。

\begin{enumerate}
\def\labelenumi{\alph{enumi})}
\item
  映射 \(J \mapsto PJ\) 是从 B 的左理想的有序集 \(\mathrm{D}(\mathrm{B}_{s})\) 到 P 的 A-子模的有序集 \(\mathrm{D}(\mathrm{P})\) 的同构。逆映射把 P 的 A-子模 M 映为 B 的左理想 \(\mathrm{J}(\mathrm{M})\)，它由 B 中使 M 包含 Pb 的元素 b 组成。
\item
  映射 \(K \mapsto KP\) 是从 A 的右理想的有序集 \(\mathrm{D}(A_{d})\) 到 P 的 B-子模的有序集 \(\mathrm{D}(P)\) 的同构。逆映射把 P 的 B-子模 V 映为 A 的右理想 \(\mathrm{K}(V)\)，它由 A 中使 V 包含 aP 的元素 a 组成。
\end{enumerate}

事实上，A-模 \(P \otimes_{B} B_{s}\) 可以典范地等同于 P。若 J 是 B 的一个左理想，则 \(P \otimes_{B} J\) 在 \(P \otimes_{B} B_{s}\) 中的典范像在此等同下对应于 PJ。于是映射 \(J \mapsto PJ\) 是从 \(D(B_{s})\) 到 \(D(P)\) 的有序集同构。设 \(J \in D(B_{s})\)。用 \(J'\) 表示 B 中使 PJ 包含 Pb 的元素 b 之集。它是 B 的包含 J 的左理想，并且我们有 \(PJ' \subset PJ\)。由于映射 \(J \mapsto PJ\) 是有序集的同构，必然有 \(PJ' = PJ\) 及 \(J = J'\)。这就证明了 a)。

断言 b) 由把断言 a) 应用于可逆 \((\mathrm{B}^{\circ},\mathrm{A}^{\circ})_{k}\)-双模 P 得出。

注意，环 A 是体当且仅当 B-模 P 是单模。

\begin{enumerate}
\def\labelenumi{\arabic{enumi})}
\setcounter{enumi}{1}
\tightlist
\item
  用 \(\mathcal{B}_{\mathrm{A}}\)、\(\mathcal{B}_{\mathrm{B}}\) 与 \(\mathcal{B}_{\mathrm{P}}\) 分别表示 A 的双边理想之集、B 的双边理想之集与 P 的 \((\mathrm{A},\mathrm{B})_k\)-子双模之集。
\end{enumerate}

\begin{enumerate}
\def\labelenumi{\alph{enumi})}
\item
  映射 \(\mathfrak{b} \mapsto \mathrm{P}\mathfrak{b}\) 是从 \(\mathcal{B}_{\mathrm{B}}\) 到 \(\mathcal{B}_{\mathrm{P}}\) 的有序集同构；逆同构把一个 \((\mathrm{A},\mathrm{B})_k\)-子双模 \(\mathrm{P}'\)（\(\mathrm{P}\) 的）映为 \(\mathrm{B}\) 的由元素 \(b\)（使得 \(\mathrm{P}b \subset \mathrm{P}'\)）组成的双边理想。
\item
  映射 \(a \mapsto aP\) 是从 \(B_{A}\) 到 \(B_{P}\) 的有序集同构；逆同构把 P 的 \((A, B)_{k}\)-子双模 \(P'\) 映为 A 的由使 \(aP \subset P'\) 的元素 a 组成的双边理想。
\end{enumerate}

事实上，设 J 是 B 的一个左理想，且 \(P' = PJ\)。那么 \(P'\) 是 P 的 A-子模，并且由例 1，理想 J 由 B 中使 \(Pb \subset P'\) 的元素 b 组成。此外，\(P'\) 是 P 的 \((A, B)_{k}\)-子双模当且仅当 J 是 B 的双边理想。于是 a) 由同前所引得出。

断言 b) 由把断言 a) 应用于可逆 \((\mathrm{B}^{\circ}, \mathrm{A}^{\circ})_{k}\)-双模 P 得出。

命题 6. —— 用 \(\mathcal{B}_{\mathrm{A}}\) 与 \(\mathcal{B}_{\mathrm{B}}\) 表示 A 的双边理想之集与 B 的双边理想之集。

\begin{enumerate}
\def\labelenumi{\alph{enumi})}
\item
  存在一个有序集同构 \(f\)，从 \(\mathcal{B}_{\mathrm{A}}\) 到 \(\mathcal{B}_{\mathrm{B}}\)，它由下列性质确定：若 \(\mathfrak{a}\) 是 \(\mathrm{A}\) 的双边理想，\(\mathfrak{b}\) 是 \(\mathrm{B}\) 的双边理想，则关系 \(f(\mathfrak{a}) = \mathfrak{b}\) 等价于 \(\mathfrak{a}\mathrm{P} = \mathrm{P}\mathfrak{b}\)。
\item
  设环 A 是交换的，从而 A 可等同于 B 的中心（VIII, 第 102 页，推论 1）。同构 \(f: \mathcal{B}_{\mathrm{A}} \to \mathcal{B}_{\mathrm{B}}\) 把 A 的理想 \(\mathfrak{a}\) 映为 B 的双边理想 \(\mathrm{Ba}\)，并且我们有 \(\mathfrak{a} = \mathrm{A} \cap \mathrm{Ba}\)。
\end{enumerate}

断言 a) 由例 2 得出。

现在设 A 是交换的，并把 A 等同于 B 的中心。设 a 是 A 的一个理想。那么 Ba 是 B 的双边理想；我们有 PBa = aP，因此 \(f(a) = \text{Ba}\)。设 \(a'\) 是 A 的理想 \(A \cap Ba\)；它含于 Ba 且包含 a，故 \(Ba'\) 等于 Ba。由于 f 是双射，由此得 \(a' = a\)。

例. —— 3) 设 V 是一个左 B-模。那么上面给出的对应把 B-模 V 的零化理想映为 A-模 \(P \otimes_{B} V\) 的零化理想。事实上，用 a 表示 A-模 \(P \otimes_{B} V\) 的零化理想，用 b 表示 B-模 V 的零化理想。设 W 是 P 的 \((A, B)_{k}\)-子双模，它由使对 V 的每个元素 v 都有 \(p \otimes v = 0\) 的元素组成。我们有包含关系 \(Pb \subset W\)；反之，对每个 \(p \in W\) 与 \(q \in Q\)，我们有 \(\theta(q \otimes p)\) 属于 b。于是，\(D(B_s)\) 的元素 b 对应于 D(P) 的元素 W。同样，\(a \in D(A_d)\) 对应于 W。

\begin{enumerate}
\def\labelenumi{\arabic{enumi})}
\setcounter{enumi}{3}
\tightlist
\item
  对 A 的每个双边理想 \(\mathfrak{a}\)，把 \(\mathbf{M}_n(\mathrm{A})\) 中由系数属于 \(\mathfrak{a}\) 的矩阵组成的子集记作 \(\mathbf{M}_n(\mathfrak{a})\)。它是 \(\mathbf{M}_n(\mathrm{A})\) 的双边理想。我们有 \(\mathbf{M}_n(\mathfrak{a})\mathrm{A}^n = \mathfrak{a}^n = \mathrm{A}^n\mathfrak{a}\)。由命题 6 推出，\(\mathbf{M}_n(\mathrm{A})\) 的每个双边理想都形如 \(\mathbf{M}_n(\mathfrak{a})\)，其中 \(\mathfrak{a}\) 是 A 的双边理想。
\end{enumerate}

注. —— 我们保留上面的假设与记号，并设 \((\mathrm{B},\mathrm{A})_{k}\)-双模 Q 是 A-模 P 的对偶 \(P^{*}\)，同构 \(\lambda\) 与 \(\theta\) 是典范映射 \(\Lambda:P\otimes_{B}P^{*}\to_{s}A_{d}\) 与 \(\Theta:P^{*}\otimes_{A}P\to_{s}B_{d}\)（VIII, 第 101 页，定理 1）。由于 A-模 P 是投射的且有限生成的，对每个 A-模 M 我们有一个典范同构 \(\vartheta_{M}:P^{*}\otimes_{A}M\to\operatorname{Hom}_{A}(P,M)\)（II, §4, No.~2, 第 271 页，推论）。我们把将构造 \(M\mapsto Q\otimes_{A}M\) 换成构造 \(M\mapsto\operatorname{Hom}_{A}(P,M)\) 后对第 3、4 小节结果的改述留给读者。

